% Compile with pLaTeX via latexmk; see README.md.
\documentclass[a4paper,12pt,english,twoside,openany,dvipdfmx]{jsbook}
\usepackage{pethesis}
\usepackage{url}
\usepackage{float}
\usepackage{amsmath}

% Edit all cover information here. The institutional wording follows the
% original cover; confirm it against your programme's submission requirements.
\title{Evaluating Information Access Systems:\\A Thesis Formatting Demonstration}
\author{First Name, Last Name}
\studentid{Student ID Number}
\thesistype{A Thesis}
\department{Department of Computer Science\\and Communications Engineering}
\graduateschool{Graduate School of Fundamental Science\\and Engineering}
\university{Waseda University}
\degree{Master of Engineering}
\submissiondate{Month DD, YYYY}
\supervisor{Prof. Tetsuya Sakai}
\researchguidance{Research on Information Access}
% Only required if you enable \copyrightpage after \maketitle.
\copyrightyear{YYYY}

% The following text and numerical examples demonstrate layout only.
% They are not research findings and must be replaced before submission.
\begin{document}
\maketitle

\begin{coverabstract}
This document demonstrates the updated English master's thesis template using
an illustrative study of information access systems. It contains a research
question, a proposed evaluation procedure, example tables, a mathematical
definition, and a discussion of limitations. All numerical results are invented
for typesetting purposes. The journal and conference references are also
explicitly fictional examples, rather than evidence for the statements below.

The layout separates chapter-opening pages from continuation pages. Each chapter
begins on a fresh page with a centred page number and no running header. Later
pages retain alternating chapter and section headers. Chapters may begin on
an odd or even page, so no empty page is inserted merely to reach a right-hand
page. The first chapter includes two continuation pages to make both running
header positions visible. The remaining chapters contain a continuation page
for comparison. These demonstration page breaks can be removed when writing
an actual thesis.

The example also demonstrates a table of contents,
cross-references, acknowledgements, a bibliography and an appendix. Together,
these elements provide a practical document for reviewing margins, heading
hierarchy, caption placement and continuous page numbering.
\end{coverabstract}

\tableofcontents

\chapter{Introduction}
\label{sec:introduction}

\section{Background and Motivation}
Information access systems help users locate material that is relevant to a
question or task. A search interface may return a ranked list of documents,
while another interface may combine retrieved passages into a short answer.
In either case, evaluation begins with a clear account of the user need and
the evidence that would satisfy it. A system that retrieves useful material
for one type of question may perform differently when the vocabulary, topic,
or expected level of detail changes.

This illustrative thesis considers how to compare two retrieval configurations
under a shared evaluation procedure. The purpose of the example is to provide
realistic paragraph lengths and document structure for checking the template.
It does not report a completed experiment. Any statements about the example
systems describe the proposed study design, and the numerical values shown
later are invented. Keeping this distinction explicit also demonstrates how
a thesis can separate an experimental plan from observed results.

\section{Research Questions}
The example asks whether an additional reranking stage improves the usefulness
of the first few retrieved documents, and whether any improvement is consistent
across different query groups. These are separate questions. An increase in an
overall mean does not establish that every group benefits, because the mean
can combine gains on some queries with losses on others. The evaluation must
therefore retain query-level observations as well as summary values.

The main comparison uses the same collection, queries and relevance judgements
for both configurations. Only the ranking procedure changes. This restriction
makes the comparison easier to interpret, although it does not eliminate all
sources of uncertainty. Later chapters discuss judgement coverage, parameter
selection and the limits of a small evaluation collection.

% Demonstration-only break: expose an ordinary continuation page.
\newpage
\section{Scope and Contributions}
\label{sec:scope}
The intended contribution of this example study is an explicit and reproducible
comparison procedure. A useful report should explain how the input collection
was prepared, how queries were selected, which settings were held constant,
and how per-query scores were aggregated. The reader should be able to trace
a value in a results table back to the observations and choices that produced
it. A single headline score is not a substitute for this account.

Three outputs are envisaged. The first is a documented baseline configuration.
The second is a controlled variant that introduces a reranking stage. The third
is an analysis that considers both average performance and variation across
queries. These outputs are described as a plan, not as completed research
contributions. In a real thesis, this section should be rewritten to reflect
what was actually implemented and established by the evidence.

\subsection{Assumptions}
The example assumes that each query has a fixed information need and that the
available judgements are used consistently for both systems. It also assumes
that the collection remains unchanged during the experiment. These assumptions
simplify the comparison, but they limit the conclusions: an offline experiment
does not directly establish how people would use the system in a live setting.
The distinction should remain visible when reporting the results.

\subsection{Reading the Example}
The document includes both numbered chapters and unnumbered back matter.
Tables and figures have independent numbering, and references to them are
resolved automatically. For example, Chapter~\ref{sec:experiments} describes the
experimental design, while Appendix~\ref{app:details} lists the information
needed to reproduce it. These references should remain correct after material
is added or removed.

% A second continuation page exposes the other odd/even running header.
\newpage
\section{Organisation of the Thesis}
Chapter~\ref{sec:related_work} explains how to organise related work around the
decisions made in an evaluation. Chapter~\ref{sec:approach} then defines the
illustrative comparison and its aggregation rule. The experimental chapter
provides example tables, and Chapter~\ref{sec:results_and_discussions} shows how
to discuss an invented result without presenting it as an empirical finding.
The final chapters distinguish conclusions from possible extensions.

This organisation places the study design before the results. That order helps
the reader determine which decisions were part of the original procedure and
which interpretations arose afterwards. A complete thesis may need additional
chapters for implementation, datasets or user studies. The template does not
require every project to follow the particular chapter structure used here.

\section{Layout Inspection}
On this continuation page, the running header should identify the current
chapter or section, according to page parity. The footer should display the
page number in the same horizontal position as on the chapter-opening page.
The opening page itself should have no running header and no header rule.
This first chapter has three pages so that both ordinary header positions can
be inspected without changing the source.

The next chapter begins on the immediately following page. Its page number
may be even or odd; no empty page is required between chapters. The forced
page breaks within these sample chapters are solely for demonstrating page
styles. During normal writing, paragraphs flow automatically and each
\verb|\chapter| command starts a new page on its own.

\chapter{Related Work}
\label{sec:related_work}

\section{Organising the Literature}
A related-work chapter should help the reader understand why the study makes
particular methodological choices. One possible organisation separates
retrieval methods, evaluation resources and analysis procedures. Another
follows the research questions directly. The choice should make comparisons
clear rather than simply reproduce the order in which papers were read.
For each relevant study, the discussion should identify its problem, evidence,
and relationship to the present work.

The citations in this demonstration illustrate formatting only. The laboratory
website entry~\cite{RSL} demonstrates a web address in a bibliography. The journal
entry~\cite{exampleArticle} and conference entry~\cite{exampleConference} are
fictional and are not sources for the methodological discussion. In a real
thesis, replace these entries with verified publications and ensure that each
citation supports the sentence in which it appears.

\section{Comparing Experimental Settings}
Two studies may use similar ranking methods while evaluating substantially
different tasks. Differences in collection size, query selection, judgement
coverage or metric definition can affect the interpretation of their scores.
A careful comparison therefore reports the experimental setting alongside the
method. It should avoid treating values from unrelated collections as if they
were measurements on the same scale under identical conditions.

A compact comparison table can summarise these settings, but the surrounding
text still needs to explain their significance. For instance, a collection
containing short factual queries may not reveal the same behaviour as a
collection of exploratory requests. Such differences can motivate an
additional experiment without implying that either earlier study was flawed.

% Demonstration-only break; remove for natural pagination in a real thesis.
\newpage
\section{Identifying a Research Gap}
A research gap should be stated in terms that can be investigated. Saying that
a topic has received limited attention is less specific than identifying an
unresolved comparison or an untested assumption. In this example, the proposed
gap concerns whether a change in average retrieval quality is accompanied by
consistent behaviour across query groups. That statement leads to observable
quantities and a concrete analysis plan.

The scope of the claim should match the literature that was examined. A thesis
can explain its search strategy, inclusion criteria and stopping point without
claiming that every relevant publication has been found. It can also separate
a practical implementation gap from a conceptual limitation. These distinctions
help prevent an overly broad novelty claim from obscuring a useful, narrower
contribution.

\section{Position of the Present Example}
The proposed comparison uses an established baseline as a reference point and
changes a single component. The role of related work is to justify that baseline,
explain reasonable alternatives and identify possible confounding factors.
A stronger or newer method is not automatically the most informative baseline
if its data preparation or computational requirements cannot be matched.

The next chapter translates these considerations into a procedure. It defines
the inputs, the two configurations and the aggregation of per-query scores.
This transition connects the literature review with the actual experiment,
so that the reader can see how earlier observations influenced the study
design rather than encountering an unrelated list of publications.

\chapter{Proposed Methods}
\label{sec:approach}

\section{System Overview}
The example compares a baseline retrieval configuration with a variant that
reranks an initial candidate set. Both configurations receive the same query
and search the same collection. The baseline returns its initial ordering,
whereas the variant applies an additional scoring function to the candidates.
Documents outside that candidate set cannot be promoted by the reranker, which
makes candidate selection an important part of the experimental description.

The implementation should record the candidate-set size, preprocessing rules,
and treatment of tied scores. These details can change the output even when
the high-level method name remains the same. Configuration files and software
versions should be saved with the result files. The appendix provides an
example checklist for this record, rather than prescribing a particular
software framework.

\section{Aggregation of Query Scores}
Let $s_i$ denote a selected evaluation score for query $i$, with $n$ queries
included in the comparison. The mean score is
\begin{equation}
  \label{eq:mean}
  \bar{s} = \frac{1}{n}\sum_{i=1}^{n}s_i.
\end{equation}
Equation~\eqref{eq:mean} assigns equal weight to each query. A different weighting
scheme would answer a different question and should be stated explicitly.
The metric used to obtain $s_i$ must also be defined; the aggregation formula
alone does not specify how relevance, rank or unjudged documents are handled.

The example retains both the mean and the individual scores. This allows the
analysis to distinguish a small improvement shared by many queries from a
large improvement concentrated in only a few cases. Neither pattern can be
reconstructed from the mean alone.

% Demonstration-only break.
\newpage
\section{Paired Comparison}
For each query, define the difference between the variant and baseline scores
using the same relevance judgements. Pairing the observations removes ambiguity
about which queries contributed to each system's mean. If a query is excluded,
the reason should be documented and the exclusion should be applied
consistently to both configurations.

A useful result file contains the query identifier, both scores and their
difference. Additional fields may describe the query group or judgement
coverage. These annotations support later analysis, but they should not be
chosen merely because they produce a favourable comparison. Decisions about
groups and exclusions belong in the experimental design wherever possible.

\subsection{Quality Checks}
Before aggregation, check for duplicate identifiers, missing outputs and
unexpected score ranges. The comparison should fail visibly when a required
system output is absent. Replacing an absent output with a score of zero may
be appropriate in some protocols, but it must be an explicit evaluation rule,
not an accidental consequence of a file-reading routine.

\section{Reproducibility}
The final report should distinguish deterministic settings from operations
that depend on a random seed. If repeated runs are used, explain what changes
between runs and how the resulting observations are summarised. Saving a seed
without recording the software and data versions is not a complete description
of the experimental environment.

The procedure should be checked on a small example before applying it to the
full collection. Manually inspecting a few query-level calculations can reveal
mismatched identifiers or inconsistent cutoffs. This check complements, rather
than replaces, automated validation of the complete result files.

\chapter{Experiments}
\label{sec:experiments}

\section{Experimental Configuration}
This chapter illustrates how to describe an experimental setup. The values
in Table~\ref{tab:environment} are hypothetical configuration values, not a
record of an experiment performed for this template. A completed thesis should
replace them with the actual settings and explain which choices were fixed
before examining the evaluation results.

\begin{table}[htbp]
  \centering
  \caption{Illustrative configuration; no experiment was executed}
  \label{tab:environment}
  \begin{tabular}{ll}\hline
    Component & Example setting\\\hline
    Candidate set & 100 documents per query\\
    Evaluation cutoff & First 10 results\\
    Configurations & Baseline and reranked variant\\
    Aggregation & Equal weight per query\\
    Output record & One paired score per query\\\hline
  \end{tabular}
\end{table}

\section{Collection and Queries}
The collection description should identify the document source, any filtering
steps and the version used in the evaluation. Query selection deserves the
same level of care. Queries used during development should be distinguished
from those reserved for the final comparison, and any overlap should be
reported rather than left implicit.

Table~\ref{tab:datasets} uses invented group sizes to demonstrate a second table
and its cross-reference. A real report would explain how each group was defined
and why the distinction is relevant to the research question. It would also
state whether group membership was assigned before the system outputs were
examined.

\begin{table}[htbp]
  \centering
  \caption{Invented query groups for the formatting example}
  \label{tab:datasets}
  \begin{tabular}{lrr}\hline
    Group & Development queries & Evaluation queries\\\hline
    Short requests & 20 & 40\\
    Detailed requests & 20 & 40\\
    All requests & 40 & 80\\\hline
  \end{tabular}
\end{table}

% Flush example floats before the continuation-page demonstration.
\clearpage
\section{Execution Procedure}
Run both configurations with the same prepared inputs and retain their complete
ranked outputs. The evaluation stage should read those outputs without altering
the rankings. Separating retrieval from scoring makes it possible to inspect
a result file, correct an evaluation bug and recompute scores without silently
changing the retrieval configuration at the same time.

Each run should have an identifier that connects its configuration, output and
log. A run log can record completion status and unexpected conditions, while
the configuration file records intended settings. If a run fails or produces
an incomplete output, resolve the issue before presenting a comparison as a
complete experiment. The final report should describe consequential deviations
from the original procedure.

\section{Analysis Plan}
The first analysis compares mean query scores. The second examines the
distribution of paired differences, and the third summarises the predefined
query groups. These analyses answer related questions, but none should be used
to imply more than it measures. A group summary based on a small number of
queries may be informative while remaining highly uncertain.

Qualitative inspection can complement these summaries. Select examples using
a stated rule, such as the largest positive and negative changes, and inspect
the retrieved documents for plausible explanations. Such examples illustrate
behaviour; they do not estimate how often that behaviour occurs in the broader
query population. The next chapter uses a small invented result table to
show how these distinctions can be written clearly.

\chapter{Results and Discussions}
\label{sec:results_and_discussions}

\section{Illustrative Results}
All values in Table~\ref{tab:results} are invented for layout inspection. They
must not be interpreted as measured performance. In this numerical example,
the variant has a larger mean in both groups, but the difference is larger
for short requests. Because the groups have equal size in the example design,
the overall means are the simple averages of the two group means.

\begin{table}[htbp]
  \centering
  \caption{Invented scores used only to demonstrate results formatting}
  \label{tab:results}
  \begin{tabular}{lrrr}\hline
    Query group & Baseline & Variant & Difference\\\hline
    Short requests & 0.420 & 0.460 & +0.040\\
    Detailed requests & 0.380 & 0.390 & +0.010\\
    Overall & 0.400 & 0.425 & +0.025\\\hline
  \end{tabular}
\end{table}

A real discussion would report the uncertainty of the comparison and connect
it to the research questions. It should avoid describing a numerical increase
as a reliable improvement without the evidence needed for that statement.
The size of a change also needs practical interpretation: a difference may
be observable in an offline score without noticeably changing the usefulness
of the system for its intended users.

% Demonstration-only break: keep the image with discussion on a continuation page.
\clearpage
\section{Figure Example}
Figure~\ref{fig:logo} demonstrates a centred image with its caption below it.
The image is the existing template logo and conveys no experimental result.
In a completed thesis, replace it with a relevant figure and explain what the
reader should learn from it in the surrounding text.

\begin{figure}[H]
  \centering
  \includegraphics[width=5cm]{waseda_logo.pdf}
  \caption{Waseda University logo, included as a formatting example}
  \label{fig:logo}
\end{figure}

\section{Interpreting Differences}
An aggregate result can conceal different query-level patterns. Some queries
may benefit from a reranking stage while others lose a relevant document near
the top of the list. Inspecting the paired scores helps identify these cases,
but the interpretation should be grounded in the retrieved material rather
than inferred from the sign of the score difference alone.

The discussion should distinguish observations from possible explanations.
For example, a query might contain terminology that appears in several
unrelated contexts. A retrieved document can reveal such an ambiguity, but
that single case does not establish a general failure mode. Additional cases
or a targeted experiment would be needed to support a broader claim.

\section{Limitations}
The illustrative design has several limitations. A fixed collection may not
reflect changes in language or document availability. Incomplete relevance
judgements may leave some useful documents unrecognised by the evaluation.
A small query set may also give an unstable picture of group-level behaviour.
Each limitation affects the scope of the conclusion in a different way.

A useful limitations section explains those effects rather than merely listing
potential problems. It should identify which conclusions remain supported,
which are tentative and which would require a new experiment. Since this
chapter contains no real observations, its only conclusion is about how a
results discussion can be structured. The following chapter preserves that
boundary when summarising the example.

\chapter{Conclusion}
\label{sec:conclusion}

\section{Summary of the Example}
This demonstration has described an illustrative comparison between a baseline
retrieval configuration and a reranked variant. The proposed procedure holds
the collection and queries fixed, retains paired query scores and examines
both overall and group summaries. It is a worked example of thesis structure,
not evidence that a particular retrieval method outperforms another.

The document connects research questions to an experimental plan and then to
a discussion of possible results. That connection is useful in a completed
thesis because it lets the reader assess whether the reported evidence answers
the questions posed at the beginning. A conclusion should revisit those
questions directly instead of introducing an unrelated set of claims.

\section{Reporting Contributions}
A final thesis should distinguish a software artifact, a methodological
contribution and an empirical finding. Implementing a system establishes that
an artifact exists; evaluating it under a documented protocol provides a
different kind of evidence. Neither automatically demonstrates that the
approach generalises beyond the setting that was studied.

For this template, the concrete artifact is the formatted document itself.
It demonstrates chapter openings with visible page numbers, continuation
pages with running headers, and uninterrupted chapter transitions on either
page parity. Tables, a figure, citations and an appendix provide additional
examples for reviewing the surrounding document structure.

% Demonstration-only break.
\newpage
\section{Boundaries of the Conclusions}
A conclusion should state the conditions under which its claims hold. Those
conditions may include the collection version, query population, relevance
judgements and computational constraints. Reporting them does not weaken an
otherwise sound result; it helps the reader understand what would need to
remain comparable when applying the finding elsewhere.

The invented scores in this document support no empirical claim. Their role
is to exercise table alignment, decimal formatting and cross-references.
Similarly, the fictional references demonstrate bibliography syntax without
establishing a literature basis for the study. Both should be replaced before
the template is used to report actual research.

\section{Closing Remarks}
The final paragraph of a thesis can explain why the established contribution
matters while keeping its scope precise. It should follow from the evidence
already presented rather than rely on a promise that future work will resolve
all remaining issues. Specific possible extensions belong in the next chapter,
where their motivation and requirements can be explained separately.

For layout inspection, compare this page with the preceding chapter-opening
page. Both should display the same centred footer style. Only this continuation
page should have a running header and a thin horizontal rule above the text.
The next chapter starts on the next page, without an intervening empty leaf.

\chapter{Future Work}
\label{sec:future_work}

\section{Broader Evaluation}
One possible extension is to evaluate additional query collections with clearly
different characteristics. Such an extension would test how sensitive the
comparison is to the initial study setting. It should preserve enough of the
original protocol to make the new evidence interpretable while explaining
any changes required by the additional data.

A broader collection is not automatically a better experiment. Its relevance
judgements, sampling procedure and documentation still determine what can be
learned from it. Before expanding the study, identify the uncertainty that the
new collection is intended to reduce. This makes the extension a targeted
investigation rather than simply a larger repetition of the same procedure.

\section{User-Centred Analysis}
Another extension would examine how people interact with the retrieved results.
An offline score represents selected aspects of retrieval quality, whereas
a user study can investigate task completion, interpretation or effort under
a specified interface. The study design would need to define the tasks and
outcomes carefully so that the observations answer a focused question.

The interface, instructions and participant population would become part of
the experimental setting. These elements should be documented alongside the
retrieval configuration. A user study should not be presented as a direct
replacement for an offline comparison; the two approaches can provide evidence
about different aspects of the same system.

% Demonstration-only break.
\newpage
\section{Efficiency and Deployment}
A further direction is to examine the computational cost of the reranking
stage. Relevant measurements might include response time, memory use and the
number of candidates processed. The evaluation should state how these values
were measured and which parts of the pipeline were included. A timing result
without the corresponding execution conditions is difficult to interpret.

The eventual deployment setting may impose a different balance between ranking
quality and computational cost. A small improvement may be useful in one
application but insufficient to justify added delay in another. Establishing
that balance requires a concrete account of the task and its constraints,
rather than a universal claim about the best configuration.

\section{Prioritising the Next Step}
A practical research plan should identify which extension would be most
informative and what resources it requires. The next step could be a targeted
error analysis, an additional collection or a measurement of runtime cost.
Its value depends on which uncertainty currently limits the interpretation
of the study, not simply on how many new components it introduces.

This completes the numbered chapters of the example. Acknowledgements and
references follow as unnumbered chapters, each with a visible page number on
its opening page. The appendix then demonstrates lettered chapter numbering.
These transitions are included so that the same page-style rules can be
checked across the complete thesis structure.


\begin{acknowledgement}
This page is a formatting example for acknowledgements. In a completed thesis,
it should identify the people and organisations whose contributions supported
the work. The wording should distinguish academic advice, technical assistance,
financial support and personal encouragement where appropriate.

Replace this paragraph with your own acknowledgements. Do not retain invented
funding information or imply that an individual endorsed the conclusions of
your thesis. This unnumbered chapter is included in the table of contents and
uses the same opening-page style as the numbered chapters: a visible footer
page number, with no running header or horizontal header rule.
\end{acknowledgement}

\bibliographystyle{unsrt}
\bibliography{reference}

\appendix
\chapter{Supplementary Details}
\label{app:details}

\section{Reproduction Record}
An appendix can preserve details that are necessary for reproduction but would
interrupt the main discussion. For the illustrative study, this includes the
collection version, query identifiers, judgement files and the settings used
for each configuration. The record should connect these inputs to the result
files described in Chapter~\ref{sec:experiments}.

A concise checklist helps the reader locate the required material:
\begin{enumerate}
  \item Record the exact collection and query versions.
  \item Save preprocessing and retrieval settings for both configurations.
  \item Retain ranked outputs and paired query-level evaluation scores.
  \item Document exclusions and unexpected conditions.
  \item State the command or procedure used to reproduce each reported table.
\end{enumerate}

An appendix should remain readable on its own while making its relationship
to the main text clear. Cross-references are preferable to manually typed page
numbers because the document may change during revision. This appendix is
labelled with a letter, and its sections use the corresponding letter prefix.

% Demonstration-only break.
\newpage
\section{Checking the Compiled Document}
This continuation page provides one final example of the running header and
centred footer. The opening page of Appendix~\ref{app:details} should show its
lettered title and a page number, but no running header. Its entry in the table
of contents should point to that opening page rather than to this continuation.

When replacing the sample content, inspect the compiled document after major
structural changes. Long headings, large tables and figures close to a page
boundary can change pagination. Those changes are normal, but captions and
cross-references should remain associated with the correct content. Remove
the demonstration-only page breaks to allow the actual thesis to flow naturally.

The cover and abstract retain their existing unnumbered appearance. The table
of contents and all pages after it display page numbers. Page counting
continues from the start of the document, so the first visible number need not
be one. This example does not introduce a separate Roman-numeral sequence for
front matter or restart Arabic numbering at the first chapter.

\end{document}
